% Generated draft: requires mathematical and source review.
\documentclass[11pt]{article}
\usepackage[margin=25mm]{geometry}
\usepackage{fontspec}
\setmainfont{Latin Modern Roman}
\usepackage{amsmath,amssymb,mathtools,unicode-math}
\setmathfont{Latin Modern Math}
\usepackage{enumitem,needspace,xurl,hyperref}
\setlength{\parindent}{0pt}
\setlength{\parskip}{5pt}
\setlength{\emergencystretch}{3em}
\newcommand{\N}{\mathbb{N}}
\newcommand{\Z}{\mathbb{Z}}
\newcommand{\Q}{\mathbb{Q}}
\newcommand{\R}{\mathbb{R}}
\newcommand{\C}{\mathbb{C}}
\newcommand{\F}{\mathbb{F}}
\newcommand{\A}{\mathbb{A}}
\newcommand{\PP}{\mathbb{P}}
\newcommand{\E}{\mathbb{E}}
\newcommand{\eps}{\varepsilon}
\newcommand{\dd}{\,\mathrm{d}}
\begin{document}
\section*{UnsolvedMath ID: 20002842}
\textbf{Source problem code:} catalogue entry
\subsection{Definitions and mathematical statement}
A Serre weight for $\operatorname{GL}_n(\F_q)$ is an isomorphism class of irreducible $\overline{\F}_p$-representations of that finite group, where $q$ is a power of $p$. Write $\tau=\rho|_{I_p}^{\mathrm{ss}}$ for the semisimplified inertial representation in the workshop's setting. Conjecture 7, attributed in the notes to ``Herzig(?)'', proposes the regular-weight recipe
\[
\operatorname{Wt}(\tau)_{\mathrm{reg}}
=\mathcal R\left(\operatorname{JH}\left(\overline{V(\tau)}\right)\right).
\]
Here $V(\tau)$ is the associated Deligne--Lusztig representation in the supplied interpretation, the bar denotes reduction modulo $p$ using an appropriate integral lattice, $\operatorname{JH}$ denotes the set of Jordan--H\"older constituents, and $\mathcal R$ is the reversal/regularization reflection. Exact definitions of $V$, regularity, $\mathcal R$, the coefficient field and the relevant local/global hypotheses are not given in this extracted record; the displayed identity is retained as a recipe with these unresolved conventions.

The source gives two accompanying examples: for $n=3$, $q=p$, an extension of $\mathcal R$ to a one-to-many operation is said to produce all predicted weights; for $n=2$, $q=p^r$, the proposed extension is said to agree with Fred's conjecture, but several extensions are available and the notes explicitly say it is unclear which is correct. Its remark says the experimental data seem better than those of Ash et al.; this is a historical report, not a proved comparison.

The aid reports broad results for the $\operatorname{GL}_2$ recipe and sufficiently generic tame results in important unitary and CM automorphic settings, without claiming unrestricted $\operatorname{GL}_n/\Q$. Its separate extension-sensitivity assertion is restricted to $p>3$ and an integer $k$ with $2<k<p-1$. Let $\omega:G_{\Q_p}\to\overline{\F}_p^\times$ be the mod-$p$ cyclotomic character and let $\rho_e$ represent an extension class
\[
e\in\operatorname{Ext}^1_{G_{\Q_p}}(1,\omega^{k-1}),\qquad
0\longrightarrow\omega^{k-1}\longrightarrow\rho_e\longrightarrow1\longrightarrow0.
\]
All these representations have the same semisimplified inertia $\omega^{k-1}\oplus1$. The aid's stated semisimple envelope has the two weights
\[
\sigma_1=\operatorname{Sym}^{k-2}(\overline{\F}_p^2),\qquad
\sigma_2=\det^{k-1}\otimes\operatorname{Sym}^{p-1-k}(\overline{\F}_p^2).
\]
For its exact local Serre-weight convention, it asserts
\[
W(\rho_e)=
\begin{cases}
\{\sigma_1,\sigma_2\},&e=0,\\
\{\sigma_1\},&e\ne0.
\end{cases}
\]
Consequently the companion weight $\sigma_2$ is present precisely for the split extension. The aid states this correction is covariant under simultaneous cyclotomic/determinant twisting, so that replacing $\rho_e$ by $\rho_e\otimes\omega^a$ replaces each weight by its $\det^a$ twist. The exact local definition of $W$, and any conditions for identifying it with globally occurring automorphic weights, are absent and must be checked. This obstruction shows that a semisimple-only envelope cannot distinguish these full extension classes; it is not itself a disproof of the original semisimple-inertia recipe.
\subsection{Short English statement}
State the regular Serre-weight recipe in terms of Deligne--Lusztig reduction and regularization, retaining its missing conventions. In the supplied generic two-dimensional example, split and nonsplit extensions have the same semisimplified inertia but differ by one companion weight in their exact weight sets.
\subsection{Sources}
\begin{itemize}
\item Local UnsolvedMath registry, permanent ID $20002842$, source code \texttt{catalogue entry}; canonical \texttt{aim-representation-theory-notes.json}, \texttt{notes[113]}.
\item Michael Volpato's notes, AIM workshop \emph{$p$-adic representations, modularity, and beyond}, February 20--24, 2006, Conjecture 7 ``Herzig(?)'', its two examples and experimental-data remark: \url{https://aimath.org/WWN/padicmodularity/padicmodularity.pdf}. The supplied verification reports the notes' warning about transcription errors from typesetting during talks. The original literature field is empty.
\item The original examples refer to ``Fred's conjecture'' and the remark to Ash et al.; no complete citations or definitive identification of Fred's formulation are supplied in this record.
\item Supplied research attempt 1: \texttt{partial\_result}, status \texttt{partially\_solved}, classification \texttt{PARTIAL-PROGRESS}, rigor \texttt{proved}, novelty confidence low. Its modern-literature assessment and generic extension-sensitivity claim have no authoritative full citations here; the aid is machine-generated, not peer reviewed, and requires independent verification before citation.
\end{itemize}
\end{document}
